\subsection{Betrag und Einheitsvektor}
\label{subsec:vector_algebra_magnitude_unit_vectors}

\begin{DefinitionBox}{Betrag eines Vektors}{vector_algebra_vector_magnitude}
  Der \emph{Betrag} \EnglishTerm{magnitude} eines Vektors
  $\vec a=(a_1,\dots,a_n)^{\mathsf T}\in\R^n$ ist seine Länge
  \[
    |\vec a|
    =\sqrt{a_1^2+a_2^2+\dots+a_n^2}.
  \]
\end{DefinitionBox}

\noindent
\begin{minipage}[c]{0.52\linewidth}
  Für $\vec a=(3,4)^{\mathsf T}$ liefert der Satz des Pythagoras
  \[
    |\vec a|
    =\Cloze{5cm}{\sqrt{3^2+4^2}=5}.
  \]
  Für $\vec b=(1,1,1)^{\mathsf T}\in\R^3$ gilt
  \[
    |\vec b|
    =\Cloze{5cm}{\sqrt{1^2+1^2+1^2}=\sqrt3}.
  \]
\end{minipage}\hfill
\begin{minipage}[c]{0.44\linewidth}
  \centering
  \CachedTikzFigure{content/resources/chapters/vector_algebra/tikz/vector_magnitude.tex}
\end{minipage}

\begin{TheoremBox}{Rechenregeln für den Betrag}{vector_algebra_magnitude_laws}
  Für $\lambda\in\R$ und $\vec a,\vec b\in\R^n$ gilt
  \[
    |\vec a|\ge0,
    \qquad
    |\lambda\vec a|=|\lambda|\,|\vec a|,
    \qquad
    |\vec a+\vec b|\le|\vec a|+|\vec b|.
  \]
  Die letzte Beziehung heißt \emph{Dreiecksungleichung}
  \EnglishTerm{triangle inequality}.
\end{TheoremBox}

\begin{center}
  \CachedTikzFigure{content/resources/chapters/vector_algebra/tikz/triangle_inequality.tex}
\end{center}

\clearpage
\begin{DefinitionBox}{Einheitsvektor}{vector_algebra_unit_vector}
  Ein \emph{Einheitsvektor} \EnglishTerm{unit vector} $\vec e$ besitzt die
  Länge
  \[
    |\vec e|=1.
  \]
  Zu jedem $\vec a\ne\vec 0$ gehört der gleichgerichtete Einheitsvektor
  \[
    \boxed{\vec e_a=\frac{\vec a}{|\vec a|}}.
  \]
\end{DefinitionBox}

\begin{ExampleBox}{Einen Vektor normieren}{vector_algebra_normalize_vector}
  Für $\vec a=(1,1)^{\mathsf T}$ gilt
  \[
    |\vec a|=\sqrt2,
    \qquad
    \vec e_a
    =\frac{\vec a}{|\vec a|}
    =\Cloze{5cm}{\frac1{\sqrt2}\begin{pmatrix}1\\1\end{pmatrix}}.
  \]
  Damit zerfällt der Vektor in Länge und Richtung:
  \[
    \boxed{\vec a=|\vec a|\,\vec e_a}.
  \]
\end{ExampleBox}
